\documentclass{article}
\usepackage{rotating}
\usepackage{tikz}
\usepackage{graphicx}
\usepackage{fontspec}
\usepackage{tocloft}
\usepackage{titletoc}
\usepackage{enumitem}

\setmainfont{LEMONMILK-Bold.otf}[Path=Assets/Fonts/]

\begin{document}

%%%%%%%%%% CAPA %%%%%%%%%%
\begin{titlepage}
\begin{tikzpicture}[remember picture,overlay]
    \fill[color=black] (current page.north west) rectangle ++(5cm,-\paperheight);
    \node[anchor=south west, text=white] at (current page.south west) {\hspace{1.5cm}\begin{turn}{90}
    \resizebox{!}{35.5pt}{\hspace{1.0cm}DOCUMENTAÇÃO DO SOFTWARE}\end{turn}};
\end{tikzpicture}

\begin{flushright}
\vspace{7.5cm}

\Huge SCANFACE

\vspace{0.5cm}

\Large Autenticação e Reconhecimento Facial

\vspace{8cm}

\normalsize Setembro 15, 2026

\end{flushright}
\end{titlepage}

\newpage
%%%%%%%%%%%%%%%%%%%%%%%%%%%%%%%%%

%%%%%%%%%% SUMÁRIO %%%%%%%%%%
\renewcommand{\contentsname}{Sumário}
\tableofcontents
\newpage
%%%%%%%%%%%%%%%%%%%%%%%%%%%%%%%%%

%%%%%%%%%% CONTEÚDO %%%%%%%%%%
\setmainfont{LEMONMILK-Bold.otf}[Path=Assets/Fonts/]
\section{Descrição}
\setmainfont{LEMONMILK-Regular.otf}[Path=Assets/Fonts/]

O ScanFace é um sistema avançado de reconhecimento e verificação facial projetado para oferecer autenticação rápida, fluida e extremamente segura em fluxos de controle de acesso. Desenvolvida com base em algoritmos de Inteligência Artificial e Visão Computacional, a solução executa a detecção de faces, o alinhamento de feições, a extração de vetores biométricos e a validação de identidade em tempo real.

A plataforma foi arquitetada para operar de forma contínua e escalável, garantindo alto nível de precisão na mitigação de acessos não autorizados e na prevenção de fraudes por falsificação (como ataques de foto ou vídeo). Graças à sua capacidade de integração versátil com APIs e bancos de dados, o ScanFace simplifica o gerenciamento de credenciais biométricas, combinando alta performance de processamento com confiabilidade operacional para ambientes corporativos, acadêmicos e residenciais.
\newpage

\setmainfont{LEMONMILK-Bold.otf}[Path=Assets/Fonts/]
\section{Plano de Desenvolvimento}
\setmainfont{LEMONMILK-Regular.otf}[Path=Assets/Fonts/]

O plano de desenvolvimento do ScanFace é dividido em etapas estratégicas para assegurar a entrega progressiva e segura da plataforma:

\vspace{0.4cm}

\textbf{Fase 1: Planejamento e Arquitetura} \\
Definição dos requisitos do sistema, modelagem da arquitetura de software, estruturação do banco de dados e desenho das rotinas de integração via API.

\vspace{0.3cm}

\textbf{Fase 2: Desenvolvimento do Core Biométrico} \\
Implementação dos módulos de visão computacional para captura de vídeo, detecção de faces, alinhamento facial e extração das características biométricas com bibliotecas especializadas.

\vspace{0.3cm}

\textbf{Fase 3: Otimização e Segurança} \\
Ajuste fino dos modelos para garantia de alta acurácia, redução do tempo de resposta na verificação em tempo real e integração de mecanismos anti-spoofing (detecção de vivacidade).

\vspace{0.3cm}

\textbf{Fase 4: Testes e Implantação} \\
Execução de testes funcionais, auditoria de segurança de dados e disponibilização da aplicação para ambiente de produção e integração final.
\newpage

\setmainfont{LEMONMILK-Bold.otf}[Path=Assets/Fonts/]
\section{Público-Alvo}
\setmainfont{LEMONMILK-Regular.otf}[Path=Assets/Fonts/]

\textbf{Administradores de TI e Segurança:} Profissionais responsáveis pelo gerenciamento de identidades, permissões de acesso e integração do sistema com a infraestrutura existente.

\vspace{0.5cm}

\textbf{Usuários Finais:} Colaboradores, clientes ou visitantes que utilizam a validação facial para autenticação rápida e sem atrito em pontos de acesso.

\vspace{0.5cm}

\textbf{Desenvolvedores e Integradores:} Equipes técnicas que utilizam as APIs do ScanFace para conectar a autenticação biométrica a outros softwares e plataformas.
\newpage

\setmainfont{LEMONMILK-Bold.otf}[Path=Assets/Fonts/]
\section{Funcionalidades}
\setmainfont{LEMONMILK-Regular.otf}[Path=Assets/Fonts/]

\subsection{Detecção e Processamento Facial em Tempo Real}
Captura contínua do fluxo de vídeo via câmera com identificação instantânea de rostos no campo de visão, aplicando rotinas de pré-processamento e alinhamento para garantir alta fidelidade na leitura.

\vspace{0.3cm}

\subsection{Autenticação Biométrica e Comparação de Embeddings}
Conversão da face detectada em um vetor numérico único (embedding biométrico) e comparação com a base de dados cadastrados para autorização ou negação imediata do acesso.

\vspace{0.3cm}

\subsection{Gestão de Usuários e Registro de Acessos (Logs)}
Módulo para cadastro de novos usuários com suas respectivas assinaturas biométricas e auditoria completa de acessos com data, hora e status das validações efetuadas.
\newpage

\setmainfont{LEMONMILK-Bold.otf}[Path=Assets/Fonts/]
\section{Tecnologias}
\setmainfont{LEMONMILK-Regular.otf}[Path=Assets/Fonts/]

O ScanFace adota uma pilha tecnológica moderna e altamente segura, priorizando alta performance nativa no Windows, separação clara de responsabilidades e proteção robusta de dados biométricos e credenciais.

\vspace{0.4cm}

\textbf{Linguagens, Plataforma e Interface}
\begin{itemize}[leftmargin=*]
    \item \textbf{Linguagem C\#:} Linguagem segura, moderna e com ótima integração ao ecossistema Windows.
    \item \textbf{Plataforma .NET 10:} Fornece alto desempenho e recursos atualizados para desenvolvimento desktop e web.
    \item \textbf{Interface Gráfica (WPF e XAML):} Utilizados para criar uma aplicação desktop nativa e intuitiva para Windows.
    \item \textbf{Arquitetura (MVVM e Clean Architecture):} Separam a interface gráfica, as regras de negócio e a infraestrutura, facilitando a manutenção e a escrita de testes.
\end{itemize}

\vspace{0.3cm}

\textbf{Banco de Dados, API e Comunicação}
\begin{itemize}[leftmargin=*]
    \item \textbf{Banco de Dados (SQLite):} Banco local leve, rápido e que não exige a configuração de um servidor separado.
    \item \textbf{API (ASP.NET Core):} Utilizada para disponibilizar a sincronização opcional de forma simples e eficiente.
    \item \textbf{Comunicação (HTTP e JSON):} Utilizados para transferir os dados cifrados entre o aplicativo desktop e a API.
\end{itemize}

\vspace{0.3cm}

\textbf{Criptografia e Autenticação}
\begin{itemize}[leftmargin=*]
    \item \textbf{Criptografia (AES-256-GCM):} Protege as credenciais e dados biométricos com garantia de confidencialidade e verificação de integridade.
    \item \textbf{Derivação de Chave (Argon2id):} Transforma a senha mestra em uma chave altamente resistente a ataques de força bruta.
    \item \textbf{Autenticação Local (Windows Hello e DPAPI):} Permite o desbloqueio seguro utilizando reconhecimento facial, impressão digital ou PIN nativo do sistema.
\end{itemize}

\vspace{0.3cm}

\textbf{Testes, Containers e CI/CD}
\begin{itemize}[leftmargin=*]
    \item \textbf{Testes (xUnit e Coverlet):} Utilizados para a execução de testes automatizados e análise contínua de cobertura de código.
    \item \textbf{Containerização (Docker):} Facilitador para a execução, testes e implantação isolada da API de sincronização.
    \item \textbf{Integração Contínua (GitHub Actions):} Automatiza os processos de compilação, execução de testes e geração das versões finais do aplicativo.
\end{itemize}
\newpage

\setmainfont{LEMONMILK-Bold.otf}[Path=Assets/Fonts/]
\section{Plano de Distribuição (opcional)}
\setmainfont{LEMONMILK-Regular.otf}[Path=Assets/Fonts/]

O modelo de distribuição do ScanFace é estruturado em duas vertentes principais:

\vspace{0.4cm}

\textbf{Disponibilização via Repositório Open-Source} \\
O código-fonte principal é disponibilizado publicamente no GitHub para a comunidade técnica, permitindo auditoria, contribuições operacionais e implantação local (*on-premises*).

\vspace{0.3cm}

\textbf{Pacotes e Conteinerização} \\
Disponibilização de imagens Docker prontas para implantação rápida em servidores localizados ou infraestruturas em nuvem, simplificando a configuração de dependências de visão computacional.
\newpage

%%%%%%%%%%%%%%%%%%%%%%%%%%%%%%%%%
\end{document}